\begin{figure}[t]
\centering
\resizebox{\textwidth}{!}{%
\begin{tikzpicture}[
    font=\sffamily,
    >=Stealth,
    % ---- Node styles ----
    cpuNode/.style={
        draw=classical!80, fill=classical!18,
        rounded corners=3pt,
        minimum width=22mm, minimum height=9mm,
        inner sep=2pt, align=center,
        font=\scriptsize\bfseries, text=classical!90!black,
        line width=0.6pt
    },
    llmNode/.style={
        draw=modelnative!80, fill=modelnative!18,
        rounded corners=3pt,
        minimum width=22mm, minimum height=9mm,
        inner sep=2pt, align=center,
        font=\scriptsize\bfseries, text=modelnative!90!black,
        line width=0.6pt
    },
    futureNode/.style={
        draw=modelnative!60, fill=modelnative!8,
        rounded corners=3pt,
        minimum width=22mm, minimum height=9mm,
        inner sep=2pt, align=center,
        font=\scriptsize\bfseries\itshape, text=modelnative!70!black,
        line width=0.6pt, dashed
    },
    yearLabel/.style={
        font=\tiny\bfseries, text=timelineGray!85!black
    },
    timelineLine/.style={
        line width=1.6pt, draw=#1!50
    },
    phaseBand/.style={
        fill=#1, fill opacity=0.08, rounded corners=4pt
    },
    curvedArrow/.style={
        -{Stealth[length=2.2mm, width=1.5mm]},
        thick, draw=#1,
        shorten <=2pt, shorten >=2pt
    },
    arrowLabel/.style={
        font=\tiny\bfseries, text=#1,
        fill=white, fill opacity=0.85, text opacity=1,
        rounded corners=1.5pt, inner xsep=3pt, inner ysep=1.2pt,
        draw=#1!30, line width=0.3pt
    },
    timelineHeader/.style={
        font=\small\bfseries, text=#1,
        fill=#1!6, rounded corners=3pt, inner sep=4pt
    }
]

% ==============================================================
%  COORDINATE SYSTEM
%  x-axis: logical year position (evenly spaced for clarity)
%    CPU:  x = 0, 3, 6, 9, 12, 15, 18
%    LLM:  x = 0, 3, 6, 6, 9, 9, 12   (some share year positions)
%  y-axis:
%    CPU timeline at y = 2.8
%    LLM timeline at y = -1.2
% ==============================================================
\def\cpuY{2.8}
\def\llmY{-1.2}

% ==============================================================
%  PHASE BACKGROUND BANDS
% ==============================================================
% CPU Phase 1: Sequential / Scalar era (1945-2004)
\fill[phaseBand=classical]
    (-1.3, \cpuY+1.1) rectangle (4.7, \cpuY-0.5);
% CPU Phase 2: Parallelism & Heterogeneity (2005-2019)
\fill[phaseBand=classical]
    (4.9, \cpuY+1.1) rectangle (13.7, \cpuY-0.5);
% CPU Phase 3: Specialization (2020+)
\fill[phaseBand=classical]
    (13.9, \cpuY+1.1) rectangle (19.5, \cpuY-0.5);

% LLM Phase 1: Pretrained Foundation (2020)
\fill[phaseBand=modelnative]
    (-1.3, \llmY+0.65) rectangle (1.7, \llmY-1.05);
% LLM Phase 2: Interfaces & Bottlenecks (2023-2024)
\fill[phaseBand=modelnative]
    (1.9, \llmY+0.65) rectangle (7.7, \llmY-1.05);
% LLM Phase 3: Heterogeneous Orchestration (2025-Future)
\fill[phaseBand=modelnative]
    (7.9, \llmY+0.65) rectangle (13.5, \llmY-1.05);

% Phase labels (subtle, inside bands)
\node[font=\tiny\itshape, text=classical!45]
    at (1.7, \cpuY+0.85) {Sequential Era};
\node[font=\tiny\itshape, text=classical!45]
    at (9.3, \cpuY+0.85) {Parallelism \& Heterogeneity};
\node[font=\tiny\itshape, text=classical!45]
    at (16.7, \cpuY+0.85) {Specialization};

\node[font=\tiny\itshape, text=modelnative!45]
    at (0.2, \llmY-0.7) {Foundation};
\node[font=\tiny\itshape, text=modelnative!45]
    at (4.8, \llmY-0.7) {Interfaces \& Bottlenecks};
\node[font=\tiny\itshape, text=modelnative!45]
    at (10.7, \llmY-0.7) {Orchestration \& Specialization};

% ==============================================================
%  TIMELINE HEADERS
% ==============================================================
\node[timelineHeader=classical] at (9, \cpuY+1.55) {CPU Architecture};
\node[timelineHeader=modelnative] at (6, \llmY+1.05) {LLM Systems};

% ==============================================================
%  TIMELINE LINES
% ==============================================================
\draw[timelineLine=classical] (-1.0, \cpuY) -- (19.0, \cpuY);
\draw[timelineLine=modelnative] (-1.0, \llmY) -- (13.0, \llmY);

% ==============================================================
%  CPU TIMELINE MILESTONES  (top row)
% ==============================================================
% 1945: ENIAC
\node[cpuNode] (cpu1) at (0, \cpuY) {ENIAC};
\node[yearLabel, below=2.5mm of cpu1] {1945};

% 1978: x86 ISA
\node[cpuNode] (cpu2) at (3, \cpuY) {x86 ISA};
\node[yearLabel, below=2.5mm of cpu2] {1978};

% 2005: Dennard Scaling Collapse
\node[cpuNode] (cpu3) at (6, \cpuY) {Dennard\\[-1pt]Collapse};
\node[yearLabel, below=2.5mm of cpu3] {2005};

% 2006: Multi-core
\node[cpuNode] (cpu4) at (9, \cpuY) {Multi-core};
\node[yearLabel, below=2.5mm of cpu4] {2006};

% 2012: big.LITTLE
\node[cpuNode] (cpu5) at (12, \cpuY) {big.LITTLE};
\node[yearLabel, below=2.5mm of cpu5] {2012};

% 2016: Heterogeneous SoC
\node[cpuNode] (cpu6) at (15, \cpuY) {Hetero.\\[-1pt]SoC};
\node[yearLabel, below=2.5mm of cpu6] {2016};

% 2020+: Domain-Specific Accelerators
\node[cpuNode] (cpu7) at (18, \cpuY) {Domain\\[-1pt]Accelerators};
\node[yearLabel, below=2.5mm of cpu7] {2020+};

% ==============================================================
%  LLM TIMELINE MILESTONES  (bottom row)
% ==============================================================
% 2020: Pretrained LLMs
\node[llmNode] (llm1) at (0, \llmY) {GPT-3};
\node[yearLabel, above=2.5mm of llm1] {2020};

% 2023: Token/Tool Interface
\node[llmNode] (llm2) at (3, \llmY) {ChatGPT\\[-1pt]Plugins};
\node[yearLabel, above=2.5mm of llm2] {2023};

% 2024: Inference Energy Wall
\node[llmNode] (llm3) at (5.5, \llmY) {KV Cache\\[-1pt]Wall};
\node[yearLabel, above=2.5mm of llm3] {2024};

% 2024: Sparse MoE
\node[llmNode] (llm4) at (8, \llmY) {Sparse\\[-1pt]MoE};
\node[yearLabel, above=2.5mm of llm4] {2024};

% 2025: Heterogeneous Model Orchestration
\node[llmNode] (llm5) at (10.5, \llmY) {Multi-Model\\[-1pt]Orchestr.};
\node[yearLabel, above=2.5mm of llm5] {2025};

% 2025: Specialized Small Models
\node[llmNode] (llm6) at (13, \llmY) {Specialized\\[-1pt]SLMs};
\node[yearLabel, above=2.5mm of llm6] {2025};

% Future: Alternative Substrates  (not connected by arrows; shown for completeness)
% (Omitted to keep the figure focused on the five structural parallels)

% ==============================================================
%  CURVED CONNECTING ARROWS  (CPU → LLM structural parallels)
% ==============================================================
% The arrows curve downward from the CPU timeline to the LLM timeline.
% We use bend left/right to create aesthetically curved paths.

% 1. Dennard Collapse → KV Cache Wall  ("Power Wall")
\draw[curvedArrow=red!70!black]
    (cpu3.south) to[bend right=18]
    node[arrowLabel=red!70!black, pos=0.45, anchor=west, xshift=1pt]
        {Power Wall}
    (llm3.north);

% 2. Multi-core → Sparse MoE  ("Parallelism Shift")
\draw[curvedArrow=blue!70!black]
    (cpu4.south) to[bend right=22]
    node[arrowLabel=blue!70!black, pos=0.45, anchor=west, xshift=1pt]
        {Parallelism Shift}
    (llm4.north);

% 3. big.LITTLE → Multi-Model Orchestration  ("Heterogeneity")
\draw[curvedArrow=orange!85!black]
    (cpu5.south) to[bend right=22]
    node[arrowLabel=orange!85!black, pos=0.45, anchor=west, xshift=1pt]
        {Heterogeneity}
    (llm5.north);

% 4. x86 ISA → ChatGPT Plugins  ("ISA Abstraction")
\draw[curvedArrow=classical!85!black]
    (cpu2.south) to[bend right=12]
    node[arrowLabel=classical!85!black, pos=0.50, anchor=east, xshift=-1pt]
        {ISA Abstraction}
    (llm2.north);

% 5. Domain Accelerators → Specialized SLMs  ("Specialization")
\draw[curvedArrow=purple!70!black]
    (cpu7.south) to[bend right=28]
    node[arrowLabel=purple!70!black, pos=0.40, anchor=west, xshift=1pt]
        {Specialization}
    (llm6.north);

\end{tikzpicture}}% end resizebox
\caption{Structural parallels between CPU and LLM system evolution.
Curved arrows connect milestones that share the same architectural driver:
Dennard scaling collapse and the inference energy wall both arise from a \emph{Power Wall};
the multi-core and MoE transitions reflect a \emph{Parallelism Shift};
big.LITTLE and heterogeneous model orchestration embody \emph{Heterogeneity};
x86 ISA and the token/tool interface represent \emph{ISA Abstraction};
and domain accelerators and specialized small models pursue \emph{Specialization}.}
\label{fig:paradigm_timeline}
\end{figure}
